\documentclass{standalone}

\usepackage[spanish]{babel} %------- Entorno de latex en español -------%
\usepackage[utf8]{inputenc} %------- Acentos nativos -------%
%------- Importar gráficas de ruta específica -------%
\usepackage{graphicx} %------- Para especificar tamaño de gráficos -------%
\graphicspath{{../img/}}
%----------------------------------------------------%
%---------- Formas para diagramas de flujo ----------%
\usepackage{tikz}
\usetikzlibrary{shapes, arrows, positioning}
\tikzset{font={\fontsize{11pt}{12}\selectfont}}
\tikzstyle{redondeado} = [rectangle, rounded corners, text width=2cm, minimum height=2em, minimum width=4em, text centered, draw=black, fill=red!30]
\tikzstyle{cuadro} = [rectangle, text width=2cm, minimum height=2em, minimum width=4em, text centered, draw=black, fill=cyan!20]
\tikzstyle{inclinado} = [trapezium, trapezium left angle=70, trapezium right angle=110, text width=1.5cm, minimum height=2em, minimum width=4em, text badly centered, draw=black, fill=green!30]
\tikzstyle{decision} = [diamond, text width=1.5cm, minimum height=2em, minimum width=4em, text badly centered, draw=black, fill=yellow!30]
\tikzstyle{burbuja} = [ellipse, text width=2cm, minimum height=2em, minimum width=4em, text badly centered, draw=black, fill=brown!30]
\tikzstyle{circulo} = [circle, text width=1mm, minimum height=1mm, minimum width=1mm, draw=black, fill=blue!30]
\tikzstyle{linea} = [draw, line width=0.2em, -latex']

\newcommand{\laptop}[1] {
  \begin{tikzpicture}[#1]
    \draw[thick, fill=black] (0, 0) -- (0.2, 0.4) -- (1.6, 0.4) -- (1.4, 0) -- cycle;
    \draw[thick, fill=white, line width=0.5em] (0.3, 0.54) -- (0.3, 1.3) -- (1.55, 1.3) -- (1.55, 0.54) -- cycle;
  \end{tikzpicture}
}

\newcommand{\bd}[1] {
  \begin{tikzpicture}[#1]
    \draw[line width=0.35em] (0, 0) arc (-45:-135:1.7em);
    \draw[line width=0.35em] (0, 0.1) arc (-45:-135:1.7em);
    \draw[line width=0.35em] (0, 0.3) arc (-45:-135:1.7em);
    \draw[line width=0.35em] (0, 0.4) arc (-45:-135:1.7em);
    \draw[line width=0.35em] (0, 0.6) arc (-45:-135:1.7em);
    \draw[line width=0.35em] (0, 0.7) arc (-45:-135:1.7em);
    \draw[line width=0.35em] (0, 0.7) arc (45:135:1.7em);
  \end{tikzpicture}
}

\newcommand{\zfm}[1] {
  \begin{tikzpicture}[#1]
    \draw[thick, fill=black] (0, 0) -- (0, 0.5) -- (0.15, 0.85) -- (0.15, 0.15) -- cycle;
    \draw[thick, fill=black] (0.05, 0) -- (0.18, 0.13) -- (0.55, 0.13) -- (0.42, 0.0) -- cycle;
    \draw[thick, fill=black] (0.19, 0.17) -- (0.19, 0.9) -- (0.6, 0.9) -- (0.6, 0.17) -- cycle;
    \draw[thick, fill=white] (0.28, 0.5) -- (0.25, 0.85) -- (0.52, 0.85) -- (0.55, 0.5) -- cycle;
  \end{tikzpicture}
}

%----------------------------------------------------%

\begin{document}

  \begin{tikzpicture}[auto, node distance = 1cm]
    \node [redondeado] (A) {Inicio};
    \node [cuadro, right = of A] (B) {Ordenar envío de huella};
    \node [cuadro, right = of B] (C) {Guardar operación en la base de datos};
    \node [cuadro, right = of C] (D) {Publicar comando modo escucha};
    \node [cuadro, below = of D] (E) {Cambiar a modo escucha};
    \node [cuadro, below = of E] (F) {Extraer huella de la base de datos};
    \node [cuadro, left = of F] (G) {Publicar ID de huella por MQTT};
    \node [cuadro, left = of G] (H) {Publicar huella por MQTT};
    \node [cuadro, left = of H] (I) {Recibir huella y enviarla al sensor};
    \node [cuadro, below = of I] (J) {Grabar huella en el sensor};
    \node [cuadro, below = of J] (K) {Publicar respuesta del sensor};
    \node [cuadro, right = of K] (L) {Guardar respuestas en la base de datos};
    \node [cuadro, right = of L] (M) {Publicar comando modo búsqueda};
    \node [cuadro, right = of M] (N) {Ingresar en modo búsqueda};
    \node [redondeado, below = of N] (O) {Fin};
    \node [label = above:{Cliente}, above = 0.5em of B] (Y1) {\laptop{}};
    \node [above = 0.5em of C] (Y2) {\bd{}};
    \node [above = 0.5em of D] (Y3) {\includegraphics[width = 2em]{ingenieria_proyecto/servidor.eps}};
    \node [right = 0.5em of E] (Y4) {\includegraphics[width = 2em]{ingenieria_proyecto/arduino.eps}};
    \node [right = 0.5em of F] (Y5) {\bd{}};
    \node [above = 0.5em of G] (Y6) {\includegraphics[width = 2em]{ingenieria_proyecto/registro.eps}};
    \node [above = 0.5em of H] (Y7) {\includegraphics[width=2em]{ingenieria_proyecto/huella.eps}};
    \node [above = 0.5em of I] (Y8) {\includegraphics[width = 2em]{ingenieria_proyecto/arduino.eps}};
    \node [left = 0.5em of J] (Y9) {\zfm{}};
    \node [left = 0.5em of K] (Y10) {\includegraphics[width = 2em]{ingenieria_proyecto/arduino.eps}};
    \node [above = 0.5em of L] (Y11) {\bd{}};
    \node [above = 0.5em of M] (Y12) {\includegraphics[width = 2em]{ingenieria_proyecto/servidor.eps}};
    \node [above = 0.5em of N] (Y13) {\includegraphics[width = 2em]{ingenieria_proyecto/arduino.eps}};

    \path [linea] (A) -> (B);
    \path [linea] (B) -> (C);
    \path [linea] (C) -> (D);
    \path [linea] (D) -> (E);
    \path [linea] (E) -> (F);
    \path [linea] (F) -> (G);
    \path [linea] (G) -> (H);
    \path [linea] (H) -> (I);
    \path [linea] (I) -> (J);
    \path [linea] (J) -> (K);
    \path [linea] (K) -> (L);
    \path [linea] (L) -> (M);
    \path [linea] (M) -> (N);
    \path [linea] (N) -> (O);
  \end{tikzpicture}

\end{document}